% ID: SEL-D0022-OB-E01
% ESTADO: cotejada. Segunda pasada visual completa y borrador compilado revisados.
% Fuente: Selectividad/matematicas_aplicadas/reserva_a_EXAMEN 5 -ANDALUCÍA.pdf
% SHA-256: 668689642ca997cf817c9f9f6608d2835efe898e3162b050d7b22b51f225d363
% Página: 2; opción B; ejercicio original 1.

\textbf{(2.5 puntos)} Una empresa fabrica dos tipos de agua de colonia, A y B. La colonia A contiene un 5\% de extracto de rosas y un 10\% de alcohol, mientras que la B se fabrica con un 10\% de extracto de rosas y un 15\% de alcohol. El precio de venta de la colonia A es de 24 \texteuro{}/litro y el de la B es de 40 \texteuro{}/litro. Se dispone de 70 litros de extracto de rosas y de 120 litros de alcohol. \textquestiondown{}Cuántos litros de cada colonia convendría fabricar para que el importe de la venta de la producción sea máximo?
